\section{Introduction}\label{sec:introduction}

A rational holomorphic line is \emph{semiample} if some positive Cartier
multiple is generated by global sections.  Such a multiple defines a
holomorphic map to projective space.  The abundance question considered
here is whether analytic nefness of a klt adjoint on a compact K\"ahler
fourfold, together with one nonzero plurisection, forces this conclusion.

Let \(\Delta\) be an effective rational Weil divisor on a normal complex
space \(X\).  We say that \(D=K_X+\Delta\) is an \emph{actual
\(\Q\)-Cartier adjoint} when an appropriate reflexive adjoint power is an
invertible holomorphic sheaf.  All multiples and sections of \(D\) then
refer to tensor powers of that line.  This convention allows \(K_X\) and
\(\Delta\) to fail to be separately \(\Q\)-Cartier.  The line is
\emph{analytically nef} if, for a fixed K\"ahler form and every
\(\varepsilon>0\), a fixed Cartier multiple has a smooth Hermitian metric
whose normalized curvature is bounded below by minus \(\varepsilon\)
times that form, using local smooth potentials on \(X\).  Section
\ref{sec:conventions} gives the full sheaf and metric conventions.
The Iitaka condition \(\kappa(X,D)\geq0\) means that a positive Cartier
multiple has a nonzero holomorphic section.

\begin{theorem}\label{thm:main}
Let \(X\) be a normal connected compact K\"ahler complex space of
dimension four.  Let \(\Delta\) be an effective rational Weil divisor such
that \((X,\Delta)\) is Kawamata log terminal and the actual adjoint
\(D=K_X+\Delta\) is \(\Q\)-Cartier.  If \(D\) is analytically nef and
\(\kappa(X,D)\geq0\), then there is an integer \(m>0\) for which
\(mD\) is Cartier and
\[
 H^0(X,\OO_X(mD))\otimes_\C\OO_X\longrightarrow\OO_X(mD)
\]
is surjective at every point.  If \(\kappa(X,D)=0\), then one may choose
\(m\) so that \(\OO_X(mD)\simeq\OO_X\).
\end{theorem}

Nonvanishing is a hypothesis.  The conclusion concerns the given
holomorphic adjoint line on \(X\), without projectivity or
\(\Q\)-factoriality of \(X\) and without a numerical-dimension
restriction.  When the Iitaka dimension is zero, the conclusion is an
actual rational-linear trivialization.

The threefold abundance theorems provide both the lower-dimensional
input and the methods behind this problem.  In the projective terminal
case, Miyaoka treated numerical dimension one
\cite{Miyaoka1988}, and Kawamata proved abundance for minimal
threefolds \cite{Kawamata1992}.  Passing to compact K\"ahler spaces
requires contraction and positivity arguments which cannot be obtained
by choosing a global ample divisor.  H\"oring--Peternell developed the
minimal-model and Mori-fiber-space theory for compact K\"ahler
threefolds \cite{HoeringPeternell2016,HoeringPeternell2015}.
Campana--H\"oring--Peternell established canonical abundance for normal
ordinary \(\Q\)-factorial compact K\"ahler threefolds with terminal
singularities \cite{CampanaHoeringPeternell2016}; the corrected
Chern-class argument is given in their appendix to
Guenancia--P\u{a}un's orbifold Bogomolov--Gieseker theorem
\cite[Appendix, Theorem~A.2]{GuenanciaPaun2025}.

For lc compact K\"ahler threefold pairs with rational boundary and nef
adjoint, Das--Ou proved semiampleness when the numerical dimension is
different from two or the Iitaka dimension is positive
\cite[Theorem~1.1]{DasOu2022}.  Their sequel treats numerical dimension
two and obtains abundance for lc compact K\"ahler threefolds
\cite[Theorem~1.2 and Corollary~1.3]{DasOu2025}.
These results generate the adjoint lines on the normal three-dimensional
components of a fourfold boundary.  They do not by themselves choose
sections agreeing on its intersections or extend those sections to the
fourfold.  Those are the two boundary problems treated here.

For positive Iitaka dimension, H\"oring--Lazi\'c--Lehn prove
semiampleness for nef klt adjoints on ordinary \(\Q\)-factorial compact
K\"ahler spaces through dimension four
\cite[Theorem~4.1]{HoeringLazicLehn2025}.  A crepant ordinary
\(\Q\)-factorial K\"ahler model projective over \(X\) places our pair
in that scope.  The sections of the pulled-back Cartier line are exactly
the sections from \(X\), so generation descends to the original line.
It remains to treat Iitaka dimension zero.

\subsection*{From a plurisection to a supported model}

Choose \(s\ne0\) in \(H^0(X,\OO_X(mD))\), and let
\[
 M=\frac1m\operatorname{div}(s)
\]
be its normalized effective rational divisor.  If \(M=0\), the section
already trivializes \(\OO_X(mD)\).  Suppose that \(M\ne0\).

The birational reduction requires projective contractions with
K\"ahler targets.  We prepare them by two constructions.  First,
Proposition~\ref{prop:finite-null} contracts the entire null locus of a
nef and big threefold class when that locus is a finite union of curves;
its initial target is normal compact analytic.  Second, conormal
direct-image vanishings and analytic thickenings extend a contraction of
a prime floor to the ambient space in
Proposition~\ref{prop:prime-floor-extension}.  These two constructions,
combined with the specified relative MMP, rationality, descent, and
positivity inputs, give the threefold contraction statement
Proposition~\ref{prop:threefold-contraction}.  It is used first on the
floors of the fourfold program and later on lower strata in the boundary
comparison argument.

The surrounding birational results are the pseudoeffective terminal
canonical threefold program of H\"oring--Peternell, the logarithmic
threefold program of Das--Hacon, and the crepant dlt models and supported
fourfold construction of Das--Hacon--P\u{a}un
\cite{HoeringPeternell2016,DasHacon2024,DasHaconPaun2024}.
The proof below states the precise imported inputs alongside the
constructions just described.

On a log resolution, we raise to one the coefficients of the strict
transforms of \(\Supp M\) and of the exceptional primes.  The resulting
adjoint has an effective representative supported on its entire reduced
floor.  The supported program then gives the model of
Proposition~\ref{prop:supported-model}: a normal ordinary
\(\Q\)-factorial compact K\"ahler dlt fourfold \((V,B)\) with effective
rational boundary and
\[
 A=K_V+B,\qquad P\sim_\Q A,\qquad S=\floor B,\qquad
 \Supp P=\Supp S,
\]
where \(A\) is an analytically nef actual \(\Q\)-Cartier adjoint,
\(P\) is an effective nonzero rational
\(\Q\)-Cartier divisor, and \(\kappa(V,A)=0\).  The equivalence is an
isomorphism of actual rational holomorphic lines.  Nonvanishing of \(P\)
follows from exceptional negativity: if its support disappeared, the
pullback of \(M\) would be a nonzero effective nef exceptional divisor.

The model also has the special projective resolution in
Lemma~\ref{lem:special-resolution}.  Its strict boundary and exceptional
support have globally smooth distinct SNC components, its exceptional
crepant coefficients are below one, and it is generically an isomorphism
on the image of every irreducible component of an intersection of distinct
strict floor components.
The first boundary result, Theorem~\ref{thm:floor-generation}, proves
that the already existing restriction \(A|_S\) is semiample for this
model.  It uses this resolution condition and requires no ambient
section.

\subsection*{Compatible sections on the whole floor}

The normal components and lower strata carry actual adjoint lines.
Lower-dimensional abundance, applied after the small models specified
below, makes these lines semiample; the threefold input is Das--Ou's
lc abundance theorem \cite[Corollary~1.3]{DasOu2025}.  We must choose
sections whose iterated residue restrictions agree through every
intersection, so that they descend to sections on the reduced space
\(S\) itself.

The restriction criterion sometimes extends a section without a
comparison.  When it requires a comparison, a perturbed program on the
lower stratum reduces the condition, in a sufficiently divisible common
degree, to the two coefficient-one markings on a projective-line fiber
of a Mori contraction.  They give a birational residue comparison,
called a link.  On a common projective graph, a link identifies the
compatible invertible meromorphic adjoint subsheaves; this equality
transports the actual residue sections and their products.  The two
markings can belong to distinct primes or to one degree-two branch, whose
normal finite Stein space carries the exchange involution.

Fujino's induction by pre-admissible and admissible sections and finite
symmetrization is the ancestry of this compatibility construction
\cite[Sections~2--4]{Fujino2000}.  Koll\'ar's sources and links describe
the related algebraic configuration with two disjoint distinguished
sections of a projective-line fiber
\cite[Definition~9, Theorem~10, and Proposition~14]{Kollar2013Sources};
the last proposition records the compatibility of even iterated residues
along these links.

For nonprojective surface strata, the links carry an additional geometric
constraint.  Restrictions of one K\"ahler form on \(V\) induce
positive-square classes on their smooth minimal surface models.  A link
from a three-dimensional stratum transports these classes modulo the
real span of divisor classes.  A returning chain begins and ends on the
same surface stratum.  After a common Cartier degree \(q\) is chosen,
the image of all returning chains on the space of \(qk\)-plurisections
is finite for each fixed integer \(k>0\).  The groupoid may have
infinitely many arrows: its surface arrows are generated only by these
links and their inverses, while point and curve comparisons are treated
separately.  Finite norm products impose the required invariance, and
finite hyperplane avoidance chooses sections nonzero at prescribed
points.  Their residue agreement through all intersections gives
generators on the whole floor.

This common-class constraint is essential.  Nonprojective K3 surfaces
can have automorphisms acting by a non-root-of-unity scalar on the
holomorphic two-form \cite[Theorems~3.5 and~4.1]{McMullen2002}.
Swapnajit Das's recent preprint states abundance for compact K\"ahler
semi-log-canonical threefolds \cite[Theorem~1.3]{Das2026}.  Its
Lemmas~7.2--7.3 and Corollary~7.4 contain a closely related
positive-class and ruling mechanism.  We give the explicit residue
construction here, including the actual meromorphic adjoint equality,
the normalized degree-two branch, and the uniform fixed-degree bound for
the link action.  The cited surface mechanism is prior art; the stated
semi-log-canonical theorem is not an input to this proof.

\subsection*{Lifting from the supported boundary}

The second boundary result is a dimension-free statement for an arbitrary
standard dlt pair.  The special resolution used in the fourfold
floor-generation application is not part of its data.

\begin{theorem}[Supported lifting]\label{thm:supported-lifting}
Let \((V,B)\) be a normal irreducible compact K\"ahler dlt pair of positive
dimension, with effective rational boundary and actual \(\Q\)-Cartier
adjoint \(A=K_V+B\).  Suppose that an effective nonzero rational
\(\Q\)-Cartier divisor \(P\) satisfies
\[
 P\sim_\Q A,\qquad \Supp P=\Supp\floor B.
\]
If the actual restriction \(A|_S\) to the whole reduced subspace
\(S=\floor B\) is semiample, then \(\kappa(V,A)\geq1\).
\end{theorem}

The support equality places the divisor of the initial section exactly
on the boundary where generation is known, and makes the pair klt away
from that divisor.  The theorem assumes generation on the entire reduced
\(S\), including its intersections; it does not assume nefness of \(A\).

Choose an effective Cartier multiple \(G=qP\) and an actual identity
\(\OO_V(G)\simeq\OO_V(qA)\) such that its restriction to \(S\) is
generated.  A generating system constructs a morphism
\(f:S\to T=\PP^b\) and an actual identity
\[
 \OO_V(G)|_S\simeq f^*\OO_T(1).
\]
On an ambient neighborhood \(W\subset V\) of each compact fiber, a
root pair and a normalized cyclic cover \(\pi:Z\to W\) replace the
supported divisor by a reduced Cartier
divisor \(E\).  Put \(I=\OO_Z(-E)\), and let \(g:E\to U\) be the
induced map over a parameter neighborhood \(U\subset T\).  The finite
lifting target is that, for every integer \(j\) and every \(k\geq1\),
\[
 g_*(I^j/I^{j+k+1})\longrightarrow g_*(I^j/I^{j+k})
\]
is an epimorphism of sheaves of complex vector spaces on \(U\).
The quotients are viewed on the underlying space of \(E\), and the
neighborhood used to lift a particular germ may depend on \(j\) and
\(k\).  Thus the assertion lifts section germs by one finite order using
the boundary map \(g\).

The obstruction to each order is a graded derivation with values in the
first cohomology of a normal layer.  A second root and a split residue
map embed that cohomology as a direct summand of the residue cohomology
of a resolved SNC divisor, retaining classes supported at special
parameters.  A local differentiation identity turns any nonzero value
of the derivation on a base coordinate into a map excluded by the Hodge
vanishing on a smooth projective parameter cover.  The same identity
then kills the remaining derivation.  Finally, cyclic invariants give
recurring divisorial layers on \(V\) with positive twists from
\(\OO_T(1)\).  Their sections grow while their higher cohomology and
the final loss from \(H^1(V,\OO_V)\) stay bounded.  This forces
\(\kappa(V,A)\geq1\).

The direct ancestors of this lifting method are the threefold arguments
of Miyaoka and Kawamata.  Miyaoka studies an effective pluricanonical
divisor on a minimal projective terminal threefold through neighborhood
covers and successive thickenings in the numerical-dimension-one case
\cite[Section~4]{Miyaoka1988}; he credits Reid with the suggestion to
analyze that divisor \cite[p.~220]{Miyaoka1988}.  Kawamata's alternative
proof of that case uses neighborhood roots, residues, compatible
thickenings, and a mixed Hodge complex \cite[Section~4]{Kawamata1992}.
The proof here establishes the stated compact K\"ahler dlt lifting
theorem.

Apply the two boundary results to the supported model.  Generation of
\(A|_S\) and Theorem~\ref{thm:supported-lifting} force positive Iitaka
dimension because \(P\ne0\), contradicting \(\kappa(V,A)=0\).
Therefore \(M=0\), and the original section \(s\) trivializes the
actual Cartier line \(\OO_X(mD)\).

Related manuscripts by OpenAI treat log abundance in the projective
characteristic-zero setting \cite{OpenAILogAbundance2026} and supported
boundary lifting \cite{OpenAILifting2026}.  The projective supplement
also states rational-boundary semi-log-canonical abundance over \(\C\)
and a dlt restriction result for a nef adjoint with an effective rational
representative whose support contains the floor, in every Cartier degree
\(m\geq2\) \cite[Corollaries~11.5--11.6]{OpenAILogAbundance2026}.
Those projective statements are separate context; both analytic boundary
arguments needed here are proved in this manuscript.

Section~\ref{sec:conventions} fixes the actual-line conventions and a
connectedness lemma.  Sections~\ref{sec:finite-null} and
\ref{sec:contraction-inputs} establish the contraction preparations used
on fourfold floors and on lower strata.  Section~\ref{sec:reduction}
constructs the supported model.  Sections~\ref{sec:adjunction}--\ref{sec:gluing}
construct compatible sections and prove generation of the existing
adjoint restriction on the entire reduced floor.
Section~\ref{sec:covers} constructs root neighborhoods and the split
residue map; Section~\ref{sec:hodge} proves the vanishing used for lifting;
and Section~\ref{sec:lifting} proves the finite-order surjections and
section growth.  Section~\ref{sec:completion} assembles the main theorem.
